\chapter{GIỚI THIỆU HỆ THỐNG}\label{chap:gioithieu}

\section{Đặt vấn đề}

ScholarLens là hệ thống truy xuất thông tin và hỏi đáp cho bài báo khoa học trên arXiv. Người dùng tìm bài bằng
câu hỏi tự nhiên (tiếng Việt hoặc tiếng Anh), thêm PDF bằng kéo-thả, và trò chuyện với chatbot. Chatbot trả lời
có trích dẫn tới đúng bài và đúng đoạn, kèm các bài liên quan tìm được trên đồ thị tri thức.

Hệ thống được xây dựng cho mini-project học phần Khai phá dữ liệu lớn và học máy, theo chủ đề tìm kiếm văn bản
khoa học. Dữ liệu thử nghiệm gồm 492 bài về truy xuất thông tin, RAG, đồ thị tri thức và mô hình ngôn ngữ.

\section{Mục tiêu}

Bảng~\ref{tab:muctieu} đối chiếu từng yêu cầu của đề bài với phần hệ thống đáp ứng và chương trình bày chi tiết.

\begin{table}[H]
\centering\small\setstretch{1.15}
\begin{tabularx}{\linewidth}{@{}P{0.28\linewidth}Lc@{}}
\toprule
\textbf{Yêu cầu của đề bài} & \textbf{ScholarLens đáp ứng bằng} & \textbf{Chương}\\
\midrule
Phát biểu bài toán, yêu cầu hệ thống & 7 tình huống sử dụng, 8 yêu cầu chức năng, 7 nhóm yêu cầu phi chức năng, yêu cầu môi trường sử dụng & 2, 3\\
Nguồn dữ liệu và cách thu thập & Crawler arXiv API (theo keyword, chuyên mục, mã bài), snapshot Kaggle, PDF người dùng tải lên & 4\\
Tổ chức hệ quản trị dữ liệu & SQLite + ChromaDB + chỉ mục BM25 + đồ thị tri thức; kiểm tra trùng ba lớp & 5\\
Hệ thống tìm kiếm đơn giản & Tìm kiếm lai BM25 + vector đa ngôn ngữ, gộp bằng RRF & 6\\
Tích hợp chatbot & Hai chế độ Chat / Research, Graph RAG, LLM Qwen2.5 tự host bằng vLLM & 7\\
Báo cáo kỹ thuật & Cấu trúc mã nguồn, triển khai, đánh giá, kiểm thử, hạn chế & 8\\
Trình bày giải pháp & Bộ slide 15 trang và kịch bản demo & 9\\
\bottomrule
\end{tabularx}
\caption{Đối chiếu yêu cầu của đề bài với hệ thống}\label{tab:muctieu}
\end{table}

\section{Kiến trúc tổng thể}

Hệ thống chạy trên hai máy (Hình~\ref{fig:kientruc}). Máy GPU (Kaggle, 2 GPU Tesla T4) làm phần tính toán nặng:
embedding hàng loạt và sinh câu trả lời bằng LLM. Máy demo chạy web app (trong Docker hoặc trực tiếp bằng Python),
giữ toàn bộ dữ liệu và phục vụ người dùng. Hai máy giao tiếp qua HTTPS bằng tunnel trycloudflare, có xác thực bằng
API key.

\begin{figure}[H]
\centering
% Sơ đồ kiến trúc tổng thể: máy GPU (Kaggle) và máy demo (Docker).
\begin{tikzpicture}[
  font=\footnotesize,
  box/.style={draw, rounded corners=2pt, align=center, text width=50mm, minimum height=11mm, fill=white,
    execute at begin node=\setstretch{1.05}},
  arr/.style={-{Stealth[length=2mm]}, thick},
  lbl/.style={font=\scriptsize, fill=white, inner sep=1pt}]
  % Máy GPU
  \node[box] (tun) at (0,0) {\textbf{cloudflared tunnel}\\https://\ldots.trycloudflare.com};
  \node[box, below=7mm of tun] (vllm) {\textbf{vLLM}: Qwen2.5-7B-Instruct-AWQ\\API tương thích OpenAI, 2 GPU};
  \node[box, below=7mm of vllm] (pipe) {\textbf{Pipeline dữ liệu (theo lô)}\\crawl $\to$ ingest $\to$ embedding $\to$ graph};
  % Máy demo
  \node[box] (web) at (8.0,0) {\textbf{Web app}\\FastAPI + giao diện chat};
  \node[box, below=7mm of web] (ret) {\textbf{Tìm kiếm lai + Graph RAG}\\BM25 + vector e5, gộp RRF};
  \node[box, below=7mm of ret] (store) {\textbf{Kho dữ liệu} \code{data/}\\SQLite, ChromaDB, BM25, đồ thị};
  \node[box, below=7mm of store] (emb) {\textbf{Embedder e5 (CPU)}\\embedding PDF mới tải lên};
  \node[box, above=9mm of web, text width=40mm, minimum height=8mm] (user) {Người dùng (trình duyệt)};
  % Khung hai máy
  \begin{scope}[on background layer]
    \node[draw, rounded corners=4pt, fill=black!4, inner sep=3.5mm, fit=(tun)(vllm)(pipe),
      label={[font=\small\bfseries]below:Máy GPU (Kaggle, 2$\times$T4)}] (gpu) {};
    \node[draw, rounded corners=4pt, fill=blue!4, inner sep=3.5mm, fit=(web)(ret)(store)(emb),
      label={[font=\small\bfseries]below:Máy demo (Docker hoặc Python)}] (demo) {};
  \end{scope}
  % Luồng
  \draw[arr] (user) -- (web);
  \draw[arr] (web) -- (ret);
  \draw[arr] (ret) -- (store);
  \draw[arr] (emb) -- (store);
  \draw[arr] (tun) -- (vllm);
  \draw[arr] (web.west) -- node[lbl, above]{HTTPS + API key} (tun.east);
  \draw[arr] (pipe.east) -- node[lbl, above]{chép \code{data/}} (store.west);
\end{tikzpicture}

\caption{Kiến trúc tổng thể của ScholarLens}\label{fig:kientruc}
\end{figure}

Câu hỏi đi từ trình duyệt vào web app, qua tìm kiếm lai và Graph RAG trên kho dữ liệu cục bộ. Chỉ bước sinh câu
trả lời đi qua tunnel tới vLLM. Pipeline dữ liệu chạy theo lô trên GPU, kết quả được chép sang máy demo một lần;
PDF thêm sau đó được embedding ngay trên máy demo.

\section{Công nghệ sử dụng}

\begin{table}[H]
\centering\small\setstretch{1.15}
\begin{tabularx}{\linewidth}{@{}P{0.25\linewidth}L@{}}
\toprule
\textbf{Tầng} & \textbf{Công nghệ}\\
\midrule
Thu thập, xử lý PDF & arXiv API~\cite{arxivapi}, PyMuPDF\\
Embedding & sentence-transformers~\cite{sbert}, \code{intfloat/multilingual-e5-base}~\cite{me5} (768 chiều, đa ngôn ngữ)\\
Lưu trữ & SQLite, ChromaDB 1.5.9 (chỉ mục HNSW~\cite{hnsw}), rank-bm25~\cite{bm25}, NetworkX~\cite{networkx}\\
LLM & vLLM~\cite{vllm} + Qwen2.5-7B-Instruct-AWQ~\cite{qwen25,awq}; có thể thay bằng Claude, Gemini, Ollama\\
Web app & FastAPI (streaming SSE), HTML/JS thuần, vis-network\\
Triển khai & Kaggle Notebook (GPU), Docker Compose, cloudflared\\
\bottomrule
\end{tabularx}
\caption{Công nghệ sử dụng theo từng tầng}\label{tab:congnghe}
\end{table}

\section{Kết quả chính}

\begin{itemize}
\item Kho 492 bài, 877 đoạn văn; đồ thị tri thức 5.598 nút và 8.258 cạnh.
\item Tìm bài theo tiêu đề bằng BM25 đạt Recall@1 = 0,957; một truy vấn tìm kiếm mất khoảng 0,3 giây trên CPU.
\item Kiểm tra trùng phát hiện 95\% tiêu đề bị biến đổi, không báo nhầm bài thật nào.
\item Trong năm bài liên quan mà đồ thị gợi ý, 37\% cùng chuyên mục với bài gốc, gấp 2,7 lần chọn ngẫu nhiên.
\end{itemize}

\section{Cấu trúc báo cáo}

Chương~\ref{chap:baitoan} phát biểu bài toán và các tình huống sử dụng. Chương~\ref{chap:yeucau} nêu yêu cầu chức
năng, phi chức năng và môi trường sử dụng. Chương~\ref{chap:dulieu} trình bày nguồn dữ liệu và cách thu thập.
Chương~\ref{chap:tochuc} mô tả tổ chức các kho dữ liệu và cơ chế kiểm tra trùng. Chương~\ref{chap:timkiem} trình
bày hệ thống tìm kiếm lai. Chương~\ref{chap:chatbot} trình bày chatbot và Graph RAG. Chương~\ref{chap:kythuat} là
báo cáo kỹ thuật: mã nguồn, triển khai, đánh giá, kiểm thử và hạn chế. Chương~\ref{chap:giaiphap} trình bày kịch
bản demo, kết luận và hướng phát triển.
